\chapter{Results and Network Analysis}
\label{chap:results}

\section{Frozen graph and corpus summary}

All results in this chapter refer to the frozen production graph
\GraphVersion{} unless a subsection explicitly names a historical or
experimental comparison. The graph contains 47,875 nodes and 225,731
relationships. It satisfies the project's structural size checks: more than
2,000 domain nodes, more than 2,000 resolved core researchers, more than 5,000
observed R1--R4 factual edges, and all ten relation types. These are structural
requirements, not semantic-accuracy evidence.

\begin{table}[!htbp]
  \centering
  \caption[Frozen graph inventory]{Node inventory and evidence classes in the frozen production graph.}
  \label{tab:graph-inventory}
  \begin{tabularx}{0.82\textwidth}{Xr}
    \toprule
    \textbf{Node label} & \textbf{Count}\\
    \midrule
    Researcher & 32,189\\
    Paper & 13,617\\
    Organization & 2,047\\
    Topic & 20\\
    Venue & 2\\
    \midrule
    \textbf{Total nodes} & \textbf{47,875}\\
    \bottomrule
  \end{tabularx}
  \vspace{0.7em}
  \begin{tabularx}{0.82\textwidth}{Xr}
    \toprule
    \textbf{Evidence class} & \textbf{Edges}\\
    \midrule
    Observed bibliographic (R1--R4) & 110,026\\
    Derived collaboration/topic (R7--R10) & 101,467\\
    Inferred topic (R5) & 14,222\\
    Curated taxonomy (R6) & 16\\
    \midrule
    \textbf{Total edges} & \textbf{225,731}\\
    \bottomrule
  \end{tabularx}
\end{table}
\FloatBarrier

The division in \cref{tab:graph-inventory} matters for the course's size
threshold: only the 110,026 observed R1--R4 edges count as factual core edges.
Derived coauthorship or inferred topic edges do not inflate that total. The
complete relation-level inventory follows in \cref{tab:relation-counts}.

\begin{table}[!htbp]
  \centering
  \caption[Accepted relationships by predicate]{Accepted relationship counts by predicate for graph \GraphVersion.}
  \label{tab:relation-counts}
  \begin{tabular}{llrr}
    \toprule
    \textbf{ID} & \textbf{Predicate} & \textbf{Evidence class} & \textbf{Edges}\\
    \midrule
    R1 & \texttt{AUTHORED} & Observed & 50,437\\
    R2 & \texttt{AFFILIATED\_WITH} & Observed & 38,664\\
    R3 & \texttt{PUBLISHED\_IN} & Observed & 13,617\\
    R4 & \texttt{CITES} & Observed & 7,308\\
    R5 & \texttt{HAS\_TOPIC} & Inferred topic & 14,222\\
    R6 & \texttt{SUBTOPIC\_OF} & Curated taxonomy & 16\\
    R7 & \texttt{COAUTHORED\_WITH} & Derived & 80,482\\
    R8 & \texttt{COLLABORATED\_WITH} & Derived & 13,031\\
    R9 & \texttt{RESEARCHES\_TOPIC} & Derived & 5,622\\
    R10 & \texttt{CONTRIBUTES\_TO\_TOPIC} & Derived & 2,332\\
    \midrule
    \multicolumn{3}{r}{\textbf{Total}} & \textbf{225,731}\\
    \bottomrule
  \end{tabular}
\end{table}

\FloatBarrier

\section{Research landscape over time}

The selected proceedings corpus grows from 2,240 papers in 2021 to 4,042 in
2025, with a dip to 1,997 in 2022. This pattern describes only the two selected
conference series and does not measure the size or growth of AI research as a
whole. AAAI contributes 9,910 papers and IJCAI 3,707 across the five-year
window.

\begin{table}[!htbp]
  \centering
  \caption[Annual corpus and collaboration counts]{Paper counts and collaboration growth by conference year. Active
  researcher counts overlap across years; first-observed coauthor pairs are
  assigned to their earliest supported paper year.}
  \label{tab:annual-output}
  \small
  \begin{tabularx}{\textwidth}{r*{5}{>{\centering\arraybackslash}X}}
    \toprule
    \textbf{Year} & \makecell{\textbf{AAAI}\\\textbf{papers}} & \makecell{\textbf{IJCAI}\\\textbf{papers}} &
    \makecell{\textbf{All}\\\textbf{papers}} & \makecell{\textbf{Active}\\\textbf{researchers}} & \makecell{\textbf{New}\\\textbf{pairs}}\\
    \midrule
    2021 & 1,654 & 586 & 2,240 & 6,337 & 12,599\\
    2022 & 1,319 & 678 & 1,997 & 6,136 & 10,730\\
    2023 & 1,578 & 639 & 2,217 & 6,941 & 12,998\\
    2024 & 2,331 & 790 & 3,121 & 9,361 & 17,447\\
    2025 & 3,028 & 1,014 & 4,042 & 13,082 & 26,708\\
    \midrule
    Total papers & 9,910 & 3,707 & 13,617 & -- & --\\
    \bottomrule
  \end{tabularx}
\end{table}

\FloatBarrier

The number of active resolved researchers and first-observed coauthor pairs
also increases in later years, especially in 2024--2025. Annual active-researcher
counts overlap because one researcher can publish in multiple years. The
``first observed pair'' count assigns a collaboration pair to the earliest
supported paper year; it is not an annual snapshot of all active relationships.

\begin{figure}[!htbp]
  \centering
  \includegraphics[width=\textwidth,height=0.73\textheight,keepaspectratio]{coverage.pdf}
  \caption[Topic-label coverage by venue and year]{Topic-label coverage for every included venue-year group. Counts
  inside each bar give labeled papers above total papers. The denominator includes all roster
  papers; coverage indicates label availability, not label correctness.}
  \label{fig:topic-coverage}
\end{figure}
\FloatBarrier

Topic labels cover 11,575 of 13,617 papers (85.00\%). The coverage chart shows
that the unlabeled share is present in every venue-year group, so topic trends
must be interpreted as trends among labeled papers. AAAI coverage changes from
85.19\% in 2021 to 84.21\% in 2025; IJCAI changes from 81.91\% to 83.43\%.
If abstained papers have a different topic distribution, labeled-paper shares
may not represent all papers.

\Needspace{23cm}
\begin{minipage}{\textwidth}
\captionsetup{type=figure}
  \begingroup\centering
  \includegraphics[width=\textwidth,height=0.73\textheight,keepaspectratio]{topic-trends.pdf}
  \caption[Fractional topic trends]{Fractional topic portfolio trends among labeled core papers from
  2021--2025. Each labeled paper contributes total mass one, divided across its
  accepted labels; AAAI and IJCAI are combined within each year. Inferred labels
  have not been human-validated. The eight largest topics by total fractional
  mass are displayed in separate panels on a shared scale.}
  \label{fig:topic-trends}
  \par\endgroup

Among labeled AAAI papers, the inferred multimodal share rises from 6.96\% to
16.42\% between 2021 and 2025. For IJCAI it rises from 7.43\% to 12.90\%.
The figures show topic shares among classifier-labeled papers rather than
independent measurements of disciplinary change. The topic test reports
substantial disagreement with its AI-assigned references, so these movements
are exploratory signals for inspection.
\end{minipage}\par

\Needspace{8\baselineskip}
\section{Researcher collaboration structure}

The researcher coauthorship projection has 32,189 resolved researchers and
80,482 distinct coauthor pairs. It contains 414 isolates, 2,936 connected
components, mean degree 5.00, and maximum degree 90. The largest component has
20,776 researchers (64.5\% of the resolved population). The remaining 2,935
components and isolates indicate that the accepted coauthorship projection is
not one connected system.

\begin{figure}[!htbp]
  \centering
  \includegraphics[width=\textwidth,height=0.73\textheight,keepaspectratio]{researcher-degree-ccdf.pdf}
  \caption[Researcher coauthor degree distribution]{Complementary cumulative distribution of distinct coauthor degree.
  All resolved researchers, including isolates, are in the denominator; both
  axes use logarithmic scales. The distribution has a long upper tail, but this
  descriptive plot does not establish a power-law model.}
  \label{fig:degree-ccdf}
\end{figure}
\FloatBarrier

The degree CCDF shows that most resolved researchers have relatively few
distinct coauthors while a smaller group has many. The long tail identifies
variation in observed collaboration breadth, not a claim about a specific
degree-distribution family. Withheld authors and papers without accepted
authorships are absent from this projection and can change its shape.

\begin{figure}[!htbp]
  \centering
  \includegraphics[width=\textwidth,height=0.73\textheight,keepaspectratio]{component-sizes.pdf}
  \caption[Researcher component sizes]{Connected component sizes in the researcher collaboration graph,
  ranked from largest to smallest. All 32,189 resolved researchers are included,
  so isolates appear as components of size one; the vertical axis is logarithmic.
 }
  \label{fig:components}
\end{figure}
\FloatBarrier

The component plot is dominated by one large component, followed by a steep
drop to small groups and isolates. This supports a qualitative description of
concentrated connectivity alongside many small collaboration groups. It does
not imply that researchers in the large component collaborate directly or share
a common research program.

\begin{figure}[!htbp]
  \centering
  \includegraphics[width=\textwidth,height=0.73\textheight,keepaspectratio]{community-topic-profiles.pdf}
  \caption[Community topic profiles]{Fractional inferred-topic profiles of selected large researcher
  communities from Louvain seed 17 and resolution 0.8. Distinct papers are
  fractionally allocated over their labels. Only stored top-five topics per
  community are shown; grey cells mean outside the stored profile, not zero
  support.}
  \label{fig:community-topics}
\end{figure}
\FloatBarrier

The community profiles show different mixtures of inferred topics among the
largest communities, alongside overlap in widely represented topics such as
representation learning. Because both community assignments and topic labels
are method-dependent, the plot is a navigation aid for inspecting groups rather
than a validated classification of research communities.

\section{Organization collaboration and topic portfolios}

The organization collaboration projection has 2,047 organizations, 13,031
accepted collaboration edges, 90 isolates, mean degree 12.73, and maximum degree
257. Its sector-mixing matrix counts distinct organization-pair by core-paper
supports and retains an explicit unknown sector. Academia--academia supports
are the largest displayed cell (10,956); academia--industry supports number
3,533. These are pair-paper supports, not unique pairs, contracts, or current
partnerships.

\begin{figure}[!htbp]
  \centering
  \includegraphics[width=\textwidth,height=0.73\textheight,keepaspectratio]{sector-mixing.pdf}
  \caption[Organization sector mixing]{Organization collaboration supports by sector pair. Counts are
  distinct accepted organization-pair by core-paper supports, shown in a
  symmetric matrix with unknown sector retained. A support requires two distinct
  authors with paper-specific organization evidence.}
  \label{fig:sector-mixing}
\end{figure}
\FloatBarrier

One prominent repeated-support candidate is Shanghai Jiao Tong University and
Tencent in inferred Representation learning, with 14 distinct core papers. The
source-linked examples include \href{https://doi.org/10.1609/aaai.v36i1.19954}{\emph{Exploiting Fine-Grained Face Forgery Clues via Progressive Enhancement Learning}}
and \href{https://doi.org/10.1609/aaai.v36i3.20193}{\emph{DIRL: Domain-Invariant Representation Learning for Generalizable Semantic Segmentation}},
both published in 2022. These records support co-publication under the graph's
affiliation rule. They do not establish a current institutional relationship or
commercial agreement. Sector labels and inferred topics remain provisional.

For inferred Representation learning, Zhejiang University has 156 supported
papers among 489 organization-associated topic-labeled papers, with a fractional
portfolio share of 25.1\%. Peking University has 139 among 381, with a 28.8\%
fractional share. Portfolio shares divide a paper across both accepted topic
labels and supported organizations, so they are not the same as raw support
fractions. These denominators omit papers without accepted institution alignment
or inferred labels.

\begin{figure}[!htbp]
  \centering
  \includegraphics[width=\textwidth,height=0.73\textheight,keepaspectratio]{organizations-network.pdf}
  \caption[Leading academia--industry collaboration pairs]{Twelve accepted
  academia--industry collaboration pairs with the greatest distinct-paper
  support across all topics; ties use stable endpoint IDs. This is an edge
  ranking from the full 2,047-organization projection, not an induced subgraph.
  Support requires distinct authors with paper-specific affiliation evidence.
  The topic-specific 14-paper example in the text uses a different filter.}
  \label{fig:organization-map}
\end{figure}
\FloatBarrier

\section{Centrality and bridge candidates}

Degree, strength, and betweenness answer different questions. Degree counts
unique partners; strength reflects the sum of fractional coauthorship weights;
betweenness reflects shortest-path brokerage under inverse-strength distances.
The configured approximate weighted betweenness ranks Bo An first at 0.042193
among 32,189 researchers. In the selected Louvain run (seed 17, resolution 0.8),
the identity belongs to a 367-member community associated with 206 distinct
papers, 180 of which have topic labels. Representation learning contributes 63
papers and 52.3333 fractional topic units, or 29.1\% of that community's
fractional topic mass.

Two source-checked works that support participation are
\href{https://doi.org/10.1609/aaai.v36i8.20851}{\emph{GearNet: Stepwise Dual Learning for Weakly Supervised Domain Adaptation}}
(2022) and \href{https://doi.org/10.1609/aaai.v37i6.25871}{\emph{Partial-Label Regression}}
(2023). They do not explain or cause the centrality score. The prior v3 bridge
leader was Dacheng Tao (0.054933); only three of the prior top twenty remain in
the corrected top twenty. On the v4 strict-evidence researcher projection,
betweenness top-20 overlap with the base projection is zero and the rank
correlation among shared researchers is 0.1882. Identity withholding changes
both the network and the population from which 256 approximation pivots are
drawn. The ranking is therefore a lead for evidence review, not an influence or
hiring score.

\begin{figure}[!htbp]
  \centering
  \includegraphics[width=\textwidth,height=0.73\textheight,keepaspectratio]{researchers-network.pdf}
  \caption[Bo An collaboration neighborhood]{Induced coauthorship subgraph on
  Bo An (OpenAlex A5017743551, gold) and 12 accepted one-hop neighbors: 13
  researchers and 24 edges. Neighbors are ranked by distinct-paper support,
  fractional weight, then stable ID. Blue edges touch the focal researcher;
  grey edges connect selected neighbors. Blue line width scales with repeated
  paper support. Layout seed is 17; positions are non-geographic. This bounded
  view is not the complete ego network or evidence for the global bridge rank.}
  \label{fig:researcher-map}
\end{figure}
\FloatBarrier

\section{Within-corpus citation structure}

The directed core citation projection contains 13,617 paper nodes, 7,308
in-scope citation edges, and 7,302 papers with no in-corpus incoming or outgoing
citation edge. Its mean total degree is 1.07 and maximum in-degree is 51. The
highest in-degree paper is \emph{Informer: Beyond Efficient Transformer for
Long Sequence Time-Series Forecasting} (AAAI 2021), with 51 in-scope citing
papers. References to papers outside the selected corpus are excluded, and
older papers have had more time to accumulate citations. The count is an
in-corpus attention measure, not a field-wide impact score.

\begin{figure}[!htbp]
  \centering
  \includegraphics[width=\textwidth,height=0.73\textheight,keepaspectratio]{top-cited-papers.pdf}
  \caption[Top ten papers by in-corpus citation count]{Highest in-degree papers in the directed core citation projection.
  The ten highest-ranked papers are shown. Each count is the number of distinct citing papers inside the release corpus;
  citations outside the corpus and unequal citation exposure time are not
  adjusted.}
  \label{fig:top-cited}
\end{figure}
\FloatBarrier

\Needspace{23cm}
\begin{minipage}{\textwidth}
\captionsetup{type=figure}
  \begingroup\centering
  \includegraphics[width=\textwidth,height=0.73\textheight,keepaspectratio]{citations-network.pdf}
  \caption[Selected incoming citations to Informer]{Selected incoming citation
  neighborhood of \emph{Informer} (AAAI 2021). Twelve of its 51 in-corpus citing
  papers are selected by descending full-projection in-degree, then stable ID.
  Only their 12 edges to the focal paper are shown; arrows point from citing to
  cited. This directed star is an edge subset, not an induced subgraph. Paper
  years accompany titles; color also distinguishes citing years. Other incoming
  citations and edges among the selected papers are omitted.}
  \label{fig:citation-map}
  \par\endgroup

\section{Findings as investigation leads}

The analysis yields three reproducible leads for further investigation. First,
the academia--industry pair/topic screen has ten candidates with at least eight
distinct supporting core papers; Shanghai Jiao Tong--Tencent, Xiamen--Tencent,
and Zhejiang--Alibaba each have 14 Representation-learning supports. Second,
Bo An ranks highest on the configured approximate betweenness measure, with a
source-traceable but policy-sensitive community context. Third, Zhejiang and
Peking universities show different fractional shares in the inferred
Representation-learning portfolio. Each finding has explicit support counts and
source examples, but remains subject to incomplete affiliations, inferred-topic
error, and researcher-identity uncertainty. The accompanying ten-person bridge
list and partner shortlist are screening devices, not recommendations.
\end{minipage}\par
